\documentclass[dvisvgm,tikz,border=3pt]{standalone}
\usepackage{amsmath,amssymb}
\usepackage{pgfplots}
\usepackage{CJKutf8}
\pgfplotsset{compat=1.18}
\usetikzlibrary{arrows.meta,calc,positioning}

\definecolor{themebg}{HTML}{2e362c}
\definecolor{themefg}{HTML}{edf4e8}
\definecolor{themegray}{HTML}{9ea897}
\definecolor{themecurve}{HTML}{7eb6ff}
\definecolor{themealert}{HTML}{ff7b7b}
\definecolor{themeamber}{HTML}{f5b942}

\pgfdeclarelayer{background}
\pgfdeclarelayer{foreground}
\pgfsetlayers{background,main,foreground}

\begin{document}
\begin{CJK*}{UTF8}{gbsn}
\pagecolor{themebg}
\begin{tikzpicture}[color=themefg, text=themefg, >=Stealth]

% ── 左侧：一般子空间和 (有重叠相交，dim(V1 ∩ V2) = 1) ──
\begin{scope}[shift={(0,0)}]
  \node[font=\bfseries\small, text=themecurve, above] at (1.5, 2.7) {一般子空间和：维数重叠};

  % 平面 V1 (晶蓝半透明)
  \draw[themecurve, fill=themecurve!20!themebg, opacity=0.7, line width=1.0pt]
    (0, 0) -- (3.6, 0) -- (4.2, 2.4) -- (0.6, 2.4) -- cycle;
  
  % 平面 V2 (暖金半透明，与 V1 相交于直线)
  \draw[themeamber, fill=themeamber!20!themebg, opacity=0.6, line width=1.0pt]
    (-0.5, -0.4) -- (3.1, -0.4) -- (2.3, 2.2) -- (-1.3, 2.2) -- cycle;

  % 交线 V1 ∩ V2 (珊瑚红高亮粗直线)
  \draw[->, themealert, line width=2.6pt] (-0.1, 0.9) -- (3.5, 0.9);

  % 基向量示意
  \draw[->, themecurve, line width=1.8pt] (1.7, 0.9) -- (2.5, 1.9);
  \draw[->, themeamber, line width=1.8pt] (1.7, 0.9) -- (1.0, 1.8);
  \fill[themefg] (1.7, 0.9) circle (2.0pt);

  % 文本标注 (带背景遮罩，消除重叠)
  \node[text=themecurve, fill=themebg, inner sep=1.5pt, font=\bfseries\small] at (4.7, 2.0) {子空间 $V_1$ ($\dim = 2$)};
  \node[text=themeamber, fill=themebg, inner sep=1.5pt, font=\bfseries\small] at (-1.2, 2.0) {子空间 $V_2$ ($\dim = 2$)};
  \node[text=themealert, fill=themebg, inner sep=1.5pt, font=\bfseries\small, below=4pt] at (2.3, 0.9) {交集 $V_1 \cap V_2$ ($\dim = 1$)};

  % 维数公式
  \node[font=\small, text=themefg] at (1.8, -1.0) {$\dim(V_1 + V_2) = 2 + 2 - 1 = 3 = \dim\mathbb{R}^3$};
\end{scope}

% ── 右侧：直和 (零交集，V1 ⊕ V2, dim(V1 ∩ V2) = 0) ──
\begin{scope}[shift={(8.4,0)}]
  \node[font=\bfseries\small, text=themeamber, above] at (1.5, 2.7) {直和：无冗余独立拼装};

  % 平面 V1 (晶蓝二维水平底面)
  \draw[themecurve, fill=themecurve!20!themebg, opacity=0.7, line width=1.0pt]
    (-0.5, 0) -- (3.1, 0) -- (3.7, 1.6) -- (0.1, 1.6) -- cycle;

  % 直线 V2 (暖金垂直穿过底面的直线)
  \draw[dashed, themeamber, line width=1.6pt] (1.6, -0.8) -- (1.6, 0.8);
  \draw[->, themeamber, line width=2.4pt] (1.6, 0.8) -- (1.6, 2.5);

  % 平面内独立向量
  \draw[->, themecurve, line width=1.8pt] (1.6, 0.8) -- (2.3, 1.4);
  \fill[themealert] (1.6, 0.8) circle (2.2pt);

  % 文本标注 (精准放置于各自独立空白区域)
  \node[text=themecurve, fill=themebg, inner sep=1.5pt, font=\bfseries\small] at (0.3, 0.3) {$V_1$ ($\dim = 2$)};
  \node[text=themeamber, fill=themebg, inner sep=1.5pt, font=\bfseries\small, right=3pt] at (1.6, 2.1) {$V_2$ ($\dim = 1$)};
  \node[text=themealert, fill=themebg, inner sep=1.5pt, font=\bfseries\small, anchor=west] at (2.6, 0.8) {唯一交点 $V_1 \cap V_2 = \{\boldsymbol{0}\}$};

  % 维数公式
  \node[font=\small, text=themeamber] at (1.8, -1.0) {$V_1 \oplus V_2 = \mathbb{R}^3 \iff \dim(V_1 + V_2) = 2 + 1 = 3$};
\end{scope}

% ── 底部总结横条 ──
\node[font=\bfseries\small, color=themecurve, anchor=north] at (6.0, -1.7) {
  核心心智模型：子空间相加是自由度合并；交集维数等于重叠自由度；直和即零重叠、表示唯一
};

\end{tikzpicture}
\end{CJK*}
\end{document}
